\documentclass[11pt]{article}
\usepackage[margin=1in]{geometry}
\usepackage{amsmath,amssymb,mathtools}
\usepackage{booktabs,longtable,array}
\usepackage{enumitem}
\usepackage{xcolor}
\usepackage{listings}
\usepackage{tikz}
\usepackage{hyperref}
\usepackage{microtype}
\usetikzlibrary{arrows.meta,positioning,fit,shapes.geometric,calc}
\hypersetup{colorlinks=true,linkcolor=black,citecolor=black,urlcolor=blue}
\lstset{basicstyle=\ttfamily\small,frame=single,columns=fullflexible,keepspaces=true,showstringspaces=false,breaklines=true}
\newcommand{\R}{\mathbb{R}}
\newcommand{\Z}{\mathbb{Z}}
\newcommand{\N}{\mathbb{N}}
\newcommand{\Qdom}{\mathcal{Q}}
\newcommand{\Pin}{\mathcal{P}_{\mathrm{in}}}
\newcommand{\Yset}{\mathcal{Y}}
\newcommand{\Hull}{\operatorname{hull}}
\newcommand{\range}{\operatorname{range}}
\newcommand{\midp}{\operatorname{mid}}
\newcommand{\rad}{\operatorname{rad}}
\newcommand{\fl}{\operatorname{fl}}
\newcommand{\down}{\operatorname{down}}
\newcommand{\up}{\operatorname{up}}
\newcommand{\abs}[1]{\left|#1\right|}
\newcommand{\norminf}[1]{\left\|#1\right\|_{\infty}}
\newcommand{\GammaE}{\Gamma}
\newcommand{\eps}{\varepsilon}
\newcommand{\epsv}{\boldsymbol{\varepsilon}}
\newcommand{\qv}{\mathbf{q}}
\newcommand{\pv}{\mathbf{p}}
\newcommand{\yv}{\mathbf{y}}
\newcommand{\xv}{\mathbf{x}}
\newcommand{\one}{\mathbf{1}}
\newcommand{\zero}{\mathbf{0}}
\title{\textbf{AutoIA V1 Mathematical Foundations and Certification Architecture}\\
A Self-Contained Specification Through Version 1.0.2}
\author{}
\date{September 2026}
\begin{document}
\maketitle

\begin{abstract}
AutoIA is a research implementation of automatic validated interval analysis for
existing engineering Python/NumPy programs.  The central objective is to accept a
deterministic engineering program, replace selected scalar and numerical operations
by validated counterparts, and return output enclosures that contain the exact real
semantics of every supported execution path for every input in a user-supplied
interval box.  The system is not based on naive line-by-line interval arithmetic
alone.  It combines explicit polynomials in the original uncertain inputs,
interval remainders, shared auxiliary error generators, semantic certification of
linear systems, path-aware execution, monotonicity and Bernstein range analysis,
outward-rounded binary64 proof kernels, and prepared repeated-query certificate
reuse.

This document is a mathematical specification of the current design through
AutoIA version 1.0.2.  Every mathematical symbol is defined before or at
its first substantive use.  The document separates \emph{proposal calculations},
which may use ordinary floating point and may be inaccurate, from \emph{proof
calculations}, whose only purpose is to establish inclusion.  In addition to the
core polynomial, interval, generator, linear-solve, branching, and prepared-query
machinery, the current V1 surface includes automatic configuration, exact
closed-endpoint comparison metadata, proof-driven domain refinement, validated
\texttt{root\_scalar} and \texttt{fsolve} semantics, fixed-data multilinear
\texttt{RegularGridInterpolator} semantics, parameter-dependent symmetric
positive-definite matrix-family semantics for log determinants and inverse
quadratic forms, parallel many-box CPU execution, optional non-trusted MPS/CUDA
proposal acceleration with independent CPU verification, and an
application-theorem interface.  The document also states current non-claims,
including the deliberate V1 fail-closed boundary at \texttt{scipy.linalg.eigh}.
\end{abstract}

\tableofcontents
\newpage

%===============================================================================
\section{Problem statement and certification target}
%===============================================================================

Let
\[
  \mathcal{P}:\R^m\rightarrow\R^r
\]
denote the mathematical semantics of a supported deterministic engineering
program.  The symbol \(m\) denotes the number of uncertain scalar input
parameters and \(r\) denotes the number of scalar outputs.  The input vector is
\[
  \pv=(p_1,\ldots,p_m)^T\in\R^m,
\]
and the output vector is
\[
  \yv=(y_1,\ldots,y_r)^T=\mathcal{P}(\pv)\in\R^r.
\]

The input-specification boundary requires distinguishing the physical interval the
user intends from the finite binary64 interval on which AutoIA executes its proof.
For uncertain input \(i\), let
\[
  I_i^\star=[a_i,b_i]\subset\R,
  \qquad a_i\le b_i,
\]
denote the \emph{intended exact-real physical interval}.  AutoIA's proof uses a
finite binary64 interval
\[
  I_i=[\underline p_i,\overline p_i]
\]
whose stored endpoints are interpreted as exact real numbers and which must
satisfy
\[
  \boxed{I_i^\star\subseteq I_i.}
\]
The intended physical box and the stored proof box are therefore
\[
  \Pin^\star=I_1^\star\times\cdots\times I_m^\star,
  \qquad
  \Pin=I_1\times\cdots\times I_m,
\]
with
\[
  \Pin^\star\subseteq\Pin.
\]
When the user directly constructs \texttt{Interval(lo, hi)} from binary64
values, the intended and proof endpoints are those stored values themselves, so
\(I_i^\star=I_i\).  When the user instead requests an exact-real enclosing
constructor, the proof box may be widened by one adjacent binary64 value at an
endpoint so that the exact intended endpoint is included.

Define the exact reachable sets over the intended and proof boxes by
\[
  \Yset_{\mathrm{intended}}^\star
  =\left\{\mathcal{P}(\pv):\pv\in\Pin^\star\right\},
  \qquad
  \Yset_{\mathrm{proof}}^\star
  =\left\{\mathcal{P}(\pv):\pv\in\Pin\right\}.
\]
Because \(\Pin^\star\subseteq\Pin\),
\[
  \Yset_{\mathrm{intended}}^\star
  \subseteq
  \Yset_{\mathrm{proof}}^\star.
\]
AutoIA certifies the stronger proof-box statement.  It does not promise to
compute the exact interval hull.  For every scalar output \(y_j\), AutoIA returns
an interval
\[
  Y_j=[\underline y_j,\overline y_j]
\]
such that
\[
  \boxed{
  y_j(\pv)\in Y_j
  \quad\text{for every }\pv\in\Pin.
  }
\]
The same interval therefore also contains every output generated by the intended
physical box \(\Pin^\star\).  For a piecewise/path-aware result, the system may
instead return a finite union of certified branch enclosures.  The reported
global lower and upper bounds are then the interval hull of that finite union.

\subsection{Three physical-input endpoint semantics}

AutoIA 1.0.2 exposes three deliberately distinct input constructors.

\paragraph{Direct binary64 endpoints.}
The existing form
\begin{lstlisting}[language=Python]
Interval(0.1, 0.3)
\end{lstlisting}
continues to mean exactly the two binary64 values that Python supplied.  No claim
is made that the source spelling \texttt{0.1} is the exact rational number
\(1/10\).  This backward-compatible interpretation is called
\texttt{binary64-exact} input semantics.

\paragraph{Exact-real enclosing endpoints.}
The form
\begin{lstlisting}[language=Python]
Interval.enclosing("0.1", "0.3")
\end{lstlisting}
interprets the strings as exact decimal real numbers and constructs finite
binary64 proof endpoints satisfying
\[
  \underline p_i\le a_i,
  \qquad
  b_i\le\overline p_i.
\]
Exact integers, \texttt{fractions.Fraction}, and \texttt{decimal.Decimal} values
are also accepted.  Ordinary Python or NumPy floating-point values are rejected
by this constructor because their pre-rounding mathematical provenance has
already been lost.

Internally, each exact decimal or rational endpoint is represented as an exact
rational number.  Conversion to binary64 is used only as a proposal.  AutoIA
compares the proposed binary64 value back to the exact rational value and, when
necessary, steps to the adjacent binary64 value in the required outward
direction.  The final inequality is then checked exactly.  If finite binary64
containment cannot be established, construction fails closed.

\paragraph{One-ulp padding of stored binary64 endpoints.}
The convenience form
\begin{lstlisting}[language=Python]
Interval.outward(0.1, 0.3)
\end{lstlisting}
moves the supplied lower stored binary64 value one adjacent representable value
toward \(-\infty\) and the upper value one adjacent representable value toward
\(+\infty\).  This is useful conservative padding, but it is \emph{not} a proof
of the provenance of an arbitrarily computed endpoint.  In particular, one ulp
cannot in general repair an unknown chain of earlier floating-point rounding
errors.

The result metadata records the selected input semantics, the supplied endpoint
specification, and the stored proof interval.  This makes the exact theorem being
certified inspectable rather than implicit.

\subsection{What is meant by program semantics}

For generic supported scalar arithmetic, stored binary64 constants in the user
program are interpreted as exact real constants equal to their stored values.
Recognized mathematical primitives are interpreted according to their documented
real mathematical semantics.  Examples include exact real \(\exp\), \(\log\),
\(\sin\), \(\cos\), \(\operatorname{erf}\), normal PDF/CDF, and the exact solution
of a recognized linear system \(A x=b\).  The semantic linear-solve primitive
therefore certifies the mathematical solution of the represented matrix family,
not the floating-point rounding behavior of NumPy's implementation of
\texttt{linalg.solve}.

Ordinary numeric constants in engineering code are intentionally \emph{not}
reinterpreted by the 1.0.2 input-specification patch.  For example, the source
literal \texttt{0.37} is first converted by Python to a binary64 value, and AutoIA
then treats that stored value as an exact real program constant.  A coefficient
whose physical uncertainty matters should be exposed as an uncertain input.  A
quantity that is genuinely intended as an exact mathematical rational can often
be written using exactly representable integers and supported exact-real
operations, for example \texttt{x/10} rather than requiring every decimal literal
in a legacy code to be replaced by an AutoIA-specific object.  Supported
mathematical primitives such as \(\exp\) are certified according to their stated
real semantics and do not rely on trusting a bare platform approximation of a
named transcendental constant.

If the user manually codes an iterative or approximate algorithm that is not a
recognized semantic primitive, the default goal is to certify the supported
program semantics as written.  AutoIA does not silently replace an arbitrary
legacy algorithm by a different mathematical problem merely because the latter
looks like the likely engineering intent.

%===============================================================================
\section{Notation used throughout the document}
%===============================================================================

This section defines the notation that appears later.

\begin{longtable}{p{0.20\textwidth}p{0.73\textwidth}}
\toprule
Symbol & Definition \\
\midrule
\endhead
\(m\) & Number of uncertain scalar physical inputs. \\
\(r\) & Number of scalar outputs. \\
\(\pv\) & Physical input vector \((p_1,\ldots,p_m)^T\). \\
\(I_i^\star\) & Intended exact-real physical interval for uncertain input \(i\). \\
\(I_i\) & Stored finite binary64 proof interval for uncertain input \(i\), with \(I_i^\star\subseteq I_i\). \\
\(\Pin^\star\) & Intended physical Cartesian input box \(\prod_i I_i^\star\). \\
\(\Pin\) & Stored binary64 proof box \(\prod_i I_i\); AutoIA certifies on this box. \\
\(p_{c,i}\) & Stored binary64 center chosen for input interval \(I_i\). \\
\(q_i\) & Centered physical coordinate \(q_i=p_i-p_{c,i}\). \\
\(h_i\) & Rigorous nonnegative half-width satisfying \(q_i\in[-h_i,h_i]\). \\
\(\qv\) & Centered physical vector \((q_1,\ldots,q_m)^T\). \\
\(\Qdom\) & Symmetric centered box \(\prod_{i=1}^m[-h_i,h_i]\). \\
\(\alpha\) & Multi-index \((\alpha_1,\ldots,\alpha_m)\) with nonnegative integer components. \\
\(|\alpha|\) & Total degree \(\alpha_1+\cdots+\alpha_m\). \\
\(\qv^\alpha\) & Monomial \(\prod_{i=1}^m q_i^{\alpha_i}\). \\
\(e_i\) & Unit multi-index whose \(i\)-th component is one and all others are zero. \\
\(P(\qv)\) & Stored point-coefficient polynomial in the original centered physical inputs. \\
\(R\) & Interval remainder, usually written \([\underline R,\overline R]\). \\
\(s\) & Number of active auxiliary generator variables. \\
\(\epsv\) & Auxiliary vector \((\eps_1,\ldots,\eps_s)^T\) with each \(\eps_j\in[-1,1]\). \\
\(G\) & Generator coefficient matrix.  For a scalar, this reduces to a row/vector of coefficients. \\
\(\Gamma_E\) & Concretization/interpretation of a validated representation \(E\). \\
\(A(\qv)\) & Parameter-dependent matrix in a semantic linear system. \\
\(b(\qv)\) & Parameter-dependent right-hand side in a semantic linear system. \\
\(X(\qv)\) & Stored polynomial predictor for a linear-system solution. \\
\(C\) & Ordinary floating-point point preconditioner used in a linear-system certificate. \\
\(B\) & Nonnegative matrix satisfying a rigorous componentwise defect bound. \\
\(c\) & Nonnegative vector satisfying a rigorous preconditioned residual bound. \\
\(\rho\) & Rigorous upper bound on \(\|B\|_\infty\). \\
\(\|B\|_\infty\) & Maximum absolute row-sum norm \(\max_i\sum_j |B_{ij}|\); since \(B\ge0\), this is \(\max_i\sum_jB_{ij}\). \\
\(\Hull(S)\) & Smallest interval containing a scalar set \(S\), or componentwise interval hull for a vector set. \\
\(\down(x)\) & A binary64 lower endpoint known to be no greater than the exact real quantity \(x\). \\
\(\up(x)\) & A binary64 upper endpoint known to be no less than the exact real quantity \(x\). \\
\bottomrule
\end{longtable}

Componentwise vector/matrix inequalities such as
\[
  |v|\le c
\]
mean \(|v_i|\le c_i\) for every component \(i\).  Similarly,
\[
  |M|\le B
\]
means \(|M_{ij}|\le B_{ij}\) for every entry.

%===============================================================================
\section{Rigorous centering of the input box}
%===============================================================================

Polynomial representations use centered coordinates.  For each user interval
\[
  I_i=[\underline p_i,\overline p_i],
\]
AutoIA chooses a finite binary64 proposal center \(p_{c,i}\).  Conceptually this
is near the arithmetic midpoint, but the proof does not require the chosen center
to be the exact midpoint.  The trusted backend then rigorously encloses
\[
  I_i-p_{c,i}
\]
using outward-rounded interval subtraction.  If the resulting interval is
\([\ell_i,u_i]\), AutoIA defines
\[
  h_i=\max(|\ell_i|,|u_i|).
\]
Therefore
\[
  I_i\subseteq[p_{c,i}-h_i,p_{c,i}+h_i],
\]
and hence every physical input in \(I_i\) admits the centered coordinate
\[
  q_i=p_i-p_{c,i}\in[-h_i,h_i].
\]
The complete centered domain is
\[
  \Qdom=\prod_{i=1}^m[-h_i,h_i].
\]

The symmetric centered domain can be slightly larger than the original user
interval because outward rounding may enlarge a boundary by one or more ulps.
This enlargement is sound but matters near function-domain boundaries.  For
example, a user interval whose exact lower endpoint is zero may produce a
centered polynomial domain whose rigorous reconstructed lower endpoint is a tiny
negative number.  AutoIA must not silently clip that negative number when
certifying \(\sqrt{x}\) or \(\log(x)\); the current implementation fails closed
in such a case unless a representation/path/domain refinement proves the required
real domain.

%===============================================================================
\section{The representation-independent enclosure invariant}
%===============================================================================

Every validated scalar value is associated with a representation object \(E\)
and a concretization map
\[
  \Gamma_E(\qv)\subseteq\R.
\]
The fundamental invariant is
\[
  \boxed{
  z(\qv)\in\Gamma_E(\qv)
  \quad\text{for all }\qv\in\Qdom,
  }
\]
where \(z\) is the exact real semantic quantity represented by \(E\).

Suppose an exact mathematical operation is denoted
\[
  z=\operatorname{op}(z_1,\ldots,z_k).
\]
A validated implementation \(\operatorname{op}^{\#}\) is sound when
\[
  \boxed{
  \Gamma_{\operatorname{op}^{\#}(E_1,\ldots,E_k)}(\qv)
  \supseteq
  \operatorname{op}\bigl(
      \Gamma_{E_1}(\qv),\ldots,\Gamma_{E_k}(\qv)
  \bigr)
  }
\]
for every \(\qv\in\Qdom\), with the right-hand side interpreted as the set of
all values of the exact operation on arguments chosen from the input sets.

Representation conversion is similarly required to be inclusion-preserving.  If
\(E_A\) is converted to \(E_B\), then
\[
  \Gamma_{E_A}(\qv)\subseteq\Gamma_{E_B}(\qv)
\]
for every \(\qv\).  Promotion priorities in the software influence strategy only;
they are not themselves part of the proof.

%===============================================================================
\section{The interval representation}
%===============================================================================

The simplest scalar representation is a single interval
\[
  E=I=[\underline z,\overline z].
\]
Its concretization is independent of \(\qv\):
\[
  \Gamma_E(\qv)=I.
\]
This representation forgets explicit dependence on the original physical inputs,
but it is robust and useful as a fallback or for calculations where interval
semantics are already natural.

The natural interval extension of an exact operation is used when supported.  For
example, if \(X=[\underline x,\overline x]\) and
\(Y=[\underline y,\overline y]\), then interval addition returns an outward
rounded enclosure of
\[
  [\underline x+\underline y,\overline x+\overline y].
\]
Interval multiplication encloses all four endpoint products.  Division requires
a denominator interval separated from zero.

%===============================================================================
\section{Point-coefficient polynomial plus interval remainder}
%===============================================================================

The principal dependency-preserving representation is
\[
  \boxed{
  z(\qv)\in P(\qv)+R,
  }
\]
where
\[
  P(\qv)=\sum_{\alpha\in\mathcal{A}} c_\alpha\qv^\alpha
\]
is a sparse polynomial with stored binary64 point coefficients \(c_\alpha\),
\(\mathcal{A}\) is a finite set of multi-indices, and
\[
  R=[\underline R,\overline R]
\]
is an interval remainder.  The concretization is
\[
  \Gamma_E(\qv)=P(\qv)+R
  =\{P(\qv)+r:r\in R\}.
\]

The stored coefficient \(c_\alpha\) is interpreted as an exact real number equal
to its binary64 value.  When a coefficient is produced by floating-point proposal
arithmetic, AutoIA separately computes a rigorous interval containing the exact
coefficient that the mathematical operation would produce, subtracts the stored
point coefficient, and ranges that coefficient error over the input domain into
\(R\).  This separation allows fast ordinary floating-point proposals while
keeping the final enclosure rigorous.

\subsection{Multi-index monomial range on a symmetric box}

For
\[
  \alpha=(\alpha_1,\ldots,\alpha_m),
\]
define
\[
  H_\alpha=\prod_{i=1}^m h_i^{\alpha_i}.
\]
The trusted implementation computes an upward-rounded upper bound for this
nonnegative product.  On the symmetric box \(\Qdom\), the monomial range is
\[
  \qv^\alpha\in
  \begin{cases}
    [0,H_\alpha], & \text{if every }\alpha_i\text{ is even},\\
    [-H_\alpha,H_\alpha], & \text{otherwise}.
  \end{cases}
\]
This parity-aware rule is sharper than blindly using \([-H_\alpha,H_\alpha]\)
for every monomial.

%===============================================================================
\section{Shared-generator polynomial representation}
%===============================================================================

The polynomial-plus-remainder representation loses correlation when the same
unstructured error is reused.  AutoIA therefore also supports
\[
  \boxed{
  z(\qv,\epsv)
  \in
  P(\qv)+\sum_{j=1}^s g_j\eps_j+R,
  \qquad \eps_j\in[-1,1].
  }
\]
For a vector of \(r\) outputs this is
\[
  \boxed{
  \mathbf{z}(\qv,\epsv)
  \in
  \mathbf{P}(\qv)+G\epsv+R,
  }
\]
where \(G\in\R^{r\times s}\) and \(R\) denotes a componentwise interval box.

The variables \(\eps_j\) are \emph{auxiliary} uncertainty variables.  They do not
represent additional physical uncertain inputs.  Their purpose is to preserve
provenance of approximation/certification error.  Two coefficients bearing the
same generator identifier refer to the same \(\eps_j\); distinct identifiers are
treated as independent coordinates in \([-1,1]^s\).

This distinction prevents an unsound shortcut.  If two unrelated calculations
happen to produce the same numerical error radius, AutoIA does not infer that
the errors are identical.  Shared provenance is created only by copying/reusing
a value or by a semantic operation that rigorously carries the same generator
through its output predictor.

\subsection{Why fan-out cancellation improves}

Suppose
\[
  y=P(\qv)+g\eps_1.
\]
If the same represented value is used twice, then
\[
  y-y
  =P-P+g\eps_1-g\eps_1
  =0.
\]
A pure interval remainder would instead behave like two independent copies of
the same interval and could fail to cancel.

\subsection{Closure under multiplication}

The representation is affine in \(\epsv\), so it is not algebraically closed
under multiplication.  If
\[
  x=P_x+G_x\epsv+R_x,
  \qquad
  y=P_y+G_y\epsv+R_y,
\]
then products such as
\[
  q_i\eps_j
  \quad\text{and}\quad
  \eps_i\eps_j
\]
appear.  The current representation keeps the constant-in-\(\qv\), first-order
in-\(\epsv\) contribution explicitly and rigorously ranges terms outside the
representation family into fresh bounded error.  Soundness is therefore retained
without pretending that \(P(\qv)+G\epsv+R\) is closed under arbitrary nonlinear
operations.

%===============================================================================
\section{Rigorous polynomial algebra}
%===============================================================================

\subsection{Addition and subtraction}

Let
\[
  x\in P_x+R_x,
  \qquad
  y\in P_y+R_y.
\]
The exact algebraic sum is
\[
  x+y\in(P_x+P_y)+(R_x+R_y).
\]
AutoIA computes point coefficient proposals for \(P_x+P_y\), rigorously encloses
coefficient-rounding errors, and adds their ranged effect to the remainder.
Subtraction is addition of the negated representation.

\subsection{Truncated multiplication}

Let the retained maximum polynomial order be \(K\).  The exact polynomial
product is
\[
  P_xP_y
  =\sum_\gamma C_\gamma\qv^\gamma,
  \qquad
  C_\gamma=\sum_{\alpha+\beta=\gamma}c_\alpha d_\beta.
\]
For every \(\gamma\) satisfying \(|\gamma|\le K\), AutoIA stores a point proposal
\(\widehat C_\gamma\).  It rigorously encloses the exact coefficient
\(C_\gamma\) and returns the difference
\[
  C_\gamma-\widehat C_\gamma
\]
as a coefficient-error interval.  That error is multiplied by a rigorous
monomial range and accumulated into the remainder.

For discarded terms with \(|\gamma|>K\), the full exact coefficient interval is
ranged over \(\Qdom\) and added to the remainder.  The user-selected order \(K\)
therefore changes sharpness and cost but not the inclusion contract.

The complete product also contains interactions with existing remainders.  These
are enclosed by ordinary interval multiplication using rigorous ranges of the
retained polynomial parts.

%===============================================================================
\section{Powers, reciprocal, and division}
%===============================================================================

\subsection{Integer powers}

For \(n\in\Z\), positive powers are constructed by validated multiplication.
For \(n<0\),
\[
  x^n=\frac{1}{x^{-n}},
\]
so AutoIA must prove that zero is excluded from the denominator range.
Negative bases are valid for integer exponents.

\subsection{Non-integer constant real powers}

For a finite constant binary64 exponent \(p\in\R\setminus\Z\), the current real
branch requires the certified base range to be strictly positive.  AutoIA may
decompose
\[
  p=n+r,
\]
where \(n\in\Z\) and \(|r|\le1/2\), then use integer powers together with
\(\sqrt{\cdot}\) for half powers or the identity
\[
  x^r=\exp(r\log x)
\]
for the remaining positive-real case.  An uncertain exponent is outside the
current supported contract.

\subsection{Reciprocal certificate}

Suppose a represented denominator \(y\) has a rigorous range
\[
  Y=[\underline Y,\overline Y]
\]
with
\[
  0\notin Y.
\]
Let \(Q\) be any stored polynomial proposal for \(1/y\).  The exact identity
\[
  \frac{1}{y}-Q
  =\frac{1-yQ}{y}
\]
allows AutoIA to bound the residual \(1-yQ\), divide by the denominator interval
\(Y\), and obtain a certified reciprocal.  The quality of \(Q\) affects only the
remainder width.

%===============================================================================
\section{Trusted binary64 interval kernel and outward rounding}
%===============================================================================

Machine rigor requires every proof-critical floating-point operation to enclose
its exact real counterpart.  AutoIA therefore distinguishes two classes of
calculation.

\begin{enumerate}[leftmargin=*]
\item \textbf{Proposal calculations.}  These may use ordinary round-to-nearest
binary64, NumPy linear algebra, heuristic scores, or approximate inverses.
Their outputs are never trusted merely because they were computed.
\item \textbf{Proof calculations.}  These establish interval inclusion and are
performed by the audited trusted backend or by operations reducible to trusted
interval primitives.
\end{enumerate}

For a scalar operation \(z=x\circ y\), the trusted backend returns binary64
endpoints
\[
  [\underline z,\overline z]
\]
known to contain the exact real value.  On the reference backend this is done by
round-to-nearest arithmetic followed by \texttt{nextafter} toward
\(-\infty\) for the lower endpoint and \(+\infty\) for the upper endpoint.

The backend construction audits the binary64 environment.  The audit checks,
among other things, the binary64 format, NumPy \texttt{float64} size,
round-to-nearest behavior, adjacent \texttt{nextafter} values, gradual
underflow/subnormals, and required vectorized arithmetic behavior.  If the
environment does not satisfy the assumptions, AutoIA fails closed.

\subsection{Why a single outward adjustment after a long reduction is insufficient}

A long dot product
\[
  s=\sum_{i=1}^n a_i b_i
\]
contains rounding error in every multiplication and every addition.  Merely
computing an ordinary BLAS result and moving it outward by one ulp does not in
general enclose the exact real dot product.

The Step-9/9b trusted fast path instead encloses every product and combines those
product intervals through a pairwise interval-addition tree.  If each leaf
contains its exact product and each internal node contains the exact sum of its
two child sets, induction proves that the root interval contains the exact dot
product.  Vectorization changes execution speed, not the mathematical proof.

The same idea is used for point-by-interval and interval-by-interval reductions,
and many independent dot-product trees may be batched simultaneously across
matrix rows/columns.

%===============================================================================
\section{Polynomial range analysis}
%===============================================================================

A validated symbolic representation is useful only if it can be converted to a
scalar interval bound.  AutoIA uses several sound range engines.

\subsection{Termwise range}

For
\[
  P(\qv)=\sum_\alpha c_\alpha\qv^\alpha,
\]
each monomial is enclosed separately and the resulting intervals are summed.
This method is inexpensive and robust, but it may lose cancellation between
terms.

\subsection{Derivative-based monotonicity}

For coordinate \(q_i\), AutoIA constructs the polynomial derivative
\[
  \frac{\partial P}{\partial q_i}
  =\sum_{\alpha:\alpha_i>0}
  \alpha_i c_\alpha\qv^{\alpha-e_i}.
\]
If a rigorous range proves
\[
  \frac{\partial P}{\partial q_i}\ge0
\]
throughout the box, then \(P\) is nondecreasing in \(q_i\); extrema occur at the
appropriate endpoint of that coordinate.  If the derivative is rigorously
nonpositive, the opposite endpoint is used.  Monotone directions can therefore
be fixed before ranging the unresolved directions.

\subsection{Bernstein enclosure}

After mapping the active box to a unit box, the polynomial can be expressed in a
multivariate Bernstein basis.  The range of a polynomial over the unit box is
contained in the interval hull of its Bernstein coefficients.  AutoIA computes
an outward-rounded version of this transformation.  Because the number of
Bernstein coefficients grows as a product of per-variable degrees, the method is
used selectively.

\subsection{Monotonicity-reduced Bernstein}

If some coordinates are already proved monotone, those coordinates are fixed at
appropriate endpoints and Bernstein conversion is applied only to unresolved
directions.  This can greatly reduce coefficient count.

\subsection{Portfolio principle}

If two independent rigorous methods produce intervals \(I_1\) and \(I_2\) that
both contain the same quantity, their intersection also contains that quantity,
provided the intersection is nonempty.  AutoIA therefore uses range-method
portfolios to improve sharpness without making correctness depend on which
method happens to be tighter.

%===============================================================================
\section{Validated elementary functions by Taylor composition}
%===============================================================================

For a scalar represented argument \(x\), AutoIA first obtains a rigorous argument
interval
\[
  X=[\underline x,\overline x].
\]
Choose a stored binary64 expansion center \(c\in X\) and write
\[
  \delta=x-c.
\]
For a smooth function \(f\), a degree-\(K\) Taylor representation is
\[
  f(x)
  =\sum_{n=0}^{K}a_n\delta^n+R_{K+1},
  \qquad
  a_n=\frac{f^{(n)}(c)}{n!}.
\]
The coefficient \(a_n\) is itself represented by a rigorous interval.  AutoIA
chooses a point proposal \(\widehat a_n\), stores it in the polynomial, and
propagates the difference between the rigorous coefficient interval and
\(\widehat a_n\) as certified error.  A rigorous analytic bound on the omitted
Taylor tail \(R_{K+1}\) is then added.

This generic composition mechanism is used for \(\exp\), \(\log\),
\(\sqrt{\cdot}\), \(\sin\), \(\cos\), \(\operatorname{erf}\), the standard normal
PDF, and the standard normal CDF, with function-specific coefficient recurrences
and tail bounds.

\subsection{Exponential}

For sufficiently small \(|r|\),
\[
  e^r=\sum_{k=0}^{N}\frac{r^k}{k!}+R_N.
\]
If \(|r|\le a<1\), one convenient rigorous tail bound is
\[
  |R_N|
  \le
  \frac{a^{N+1}}{(N+1)!(1-a)}.
\]
The trusted scalar kernel repeatedly halves a rigorous argument interval until it
is small enough for the series, then reconstructs the result by trusted squaring.

\subsection{Logarithm}

For \(x>0\), binary exponent extraction produces
\[
  x=2^k m,
\]
with a scaled mantissa \(m\).  Then
\[
  \log x=k\log2+\log m.
\]
Using
\[
  z=\frac{m-1}{m+1},
\]
we have
\[
  \log m=2\operatorname{atanh}(z),
\]
which is evaluated by a rigorously bounded odd series.  The constant \(\log2\)
is itself generated rigorously rather than trusted from a platform library.

\subsection{Sine and cosine}

A rigorous interval for \(\pi\) is generated from Machin's formula.  An ordinary
floating-point integer reduction index may be proposed, but the proof forms the
reduced argument with the rigorous \(\pi\) interval.  Taylor sine/cosine bounds
are used when the reduced interval is small enough; otherwise a safe range such
as \([-1,1]\) is available.

\subsection{Error function and normal distribution functions}

The error function is
\[
  \operatorname{erf}(x)
  =\frac{2}{\sqrt\pi}\int_0^x e^{-t^2}\,dt.
\]
Moderate positive arguments use a bounded alternating Maclaurin series.  Larger
arguments use the complementary error function and an integration-by-parts
asymptotic bracket.  Odd symmetry handles negative arguments.

The standard normal density and CDF are
\[
  \phi(x)=\frac{1}{\sqrt{2\pi}}e^{-x^2/2},
\]
\[
  \Phi(x)=\frac12\left(1+\operatorname{erf}\left(\frac{x}{\sqrt2}\right)\right).
\]
Their polynomial Taylor coefficients are generated using differential recurrences,
including
\[
  \phi'(x)+x\phi(x)=0
\]
and
\[
  \Phi'(x)=\phi(x).
\]
A normal distribution with constant mean \(\mu\) and standard deviation
\(\sigma>0\) is handled using
\[
  z=\frac{x-\mu}{\sigma},
  \qquad
  f_X(x)=\frac{\phi(z)}{\sigma},
  \qquad
  F_X(x)=\Phi(z).
\]

%===============================================================================
\section{Semantic certification of parameter-dependent linear systems}
%===============================================================================

Consider a represented family of linear systems
\[
  A(\qv)x(\qv)=b(\qv),
\]
where \(A(\qv)\in\R^{n\times n}\), \(x(\qv)\in\R^n\), and
\(b(\qv)\in\R^n\).  The matrix and right-hand side may have retained polynomial
coefficients of arbitrary supported order plus rigorous remainders.

\subsection{General polynomial solution predictor}

Write
\[
  A(\qv)\approx\sum_\beta A_\beta\qv^\beta,
  \qquad
  b(\qv)\approx\sum_\beta b_\beta\qv^\beta,
\]
and propose
\[
  X(\qv)=\sum_{|\alpha|\le K}x_\alpha\qv^\alpha.
\]
Formal coefficient matching gives the recurrence
\[
  \boxed{
  A_0x_\alpha
  =
  b_\alpha
  -
  \sum_{\substack{\beta\ne0\\\beta\le\alpha}}
      A_\beta x_{\alpha-\beta}.
  }
\]
Here \(\beta\le\alpha\) means \(\beta_i\le\alpha_i\) for every coordinate.
The recurrence is only a predictor generator.  AutoIA may use an ordinary
floating-point solve, a different reference matrix, a family-hull midpoint, or
another finite proposal.

\subsection{Residual identity for an arbitrary predictor}

Let \(X\) be whatever predictor is actually stored.  Define the true solution
error
\[
  e=x-X.
\]
Define the residual
\[
  r=b-AX.
\]
Because \(Ax=b\),
\[
  A(X+e)=b,
\]
so
\[
  \boxed{Ae=r.}
\]
This identity is exact and does not depend on how \(X\) was obtained.

Choose any finite point matrix \(C\) as a preconditioner.  Then
\[
  CAe=Cr.
\]
Add and subtract \(e\):
\[
  e-(I-CA)e=Cr,
\]
and hence
\[
  \boxed{
  e=Cr+(I-CA)e.
  }
\]

Suppose trusted calculations establish a nonnegative vector \(c\) and
nonnegative matrix \(B\) satisfying, for every admissible input,
\[
  |Cr|\le c,
\]
\[
  |I-CA|\le B.
\]
Taking componentwise absolute values gives
\[
  |e|\le c+B|e|.
\]

\subsection{Contraction theorem}

Define
\[
  \rho=\|B\|_\infty=\max_i\sum_jB_{ij}.
\]
If a rigorous upper bound proves
\[
  \boxed{\rho<1,}
\]
then the family is nonsingular throughout the represented domain and
\[
  \boxed{
  |e|\le(I-B)^{-1}c.
  }
\]
A simpler scalar consequence is obtained by defining
\[
  \eta=\|c\|_\infty,
  \qquad
  \delta=1-\rho>0.
\]
Then
\[
  \|e\|_\infty\le\frac{\eta}{\delta}.
\]
The implementation computes an outward-rounded upper bound for this ratio and
may subsequently refine the componentwise correction by the monotone iteration
\[
  d^{(k+1)}=\min\left(d^{(k)},\;\up(c+Bd^{(k)})\right),
\]
starting from a safe scalar vector bound.

\subsection{Preconditioner portfolio}

The center matrix inverse is a useful proposal but can be a poor contraction
preconditioner.  AutoIA also constructs a rigorous interval hull of the matrix
family, chooses a finite representative such as its componentwise midpoint, and
proposes the inverse of that representative.  Every candidate is checked against
the complete matrix family.  Therefore a midpoint that is not physically
realizable still yields a sound certificate if the final defect bound passes.

%===============================================================================
\section{Validated symmetric positive-definite matrix families}
%===============================================================================

Many engineering calculations use a symmetric matrix whose entries depend on
common uncertain inputs.  Let
\[
  A(\pv)\in\R^{n\times n},\qquad \pv\in\Pin,
\]
where \(n\) is the matrix dimension.  Assume the matrix is intended to satisfy
\[
  A(\pv)=A(\pv)^T
\]
for every admissible \(\pv\).  The notation
\[
  A(\pv)\succ0
\]
means that the matrix is symmetric positive definite (SPD), namely
\[
  \xv^T A(\pv)\xv>0
  \quad\text{for every nonzero }\xv\in\R^n.
\]
This structure appears in covariance matrices, structural stiffness and mass
matrices, Fisher information matrices, normal-equation matrices, Gaussian
likelihoods, Kalman filters, and many other applications.

A naive entrywise interval matrix can be much larger than the actual matrix
family because every entry is then allowed to vary independently.  AutoIA
therefore keeps the original shared physical dependence in the scalar matrix
entries and applies a structured preconditioning theorem.

\subsection{Constant proposal preconditioner}

Let \(\pv_c\) be a nominal point, usually the center of the physical input box,
and let
\[
  A_c=A(\pv_c)
\]
be an ordinary binary64 proposal matrix.  AutoIA attempts an ordinary numerical
Cholesky factorization
\[
  A_c\approx L L^T,
\]
where \(L\) is lower triangular with positive diagonal.  This factorization is
\emph{not} trusted as a proof.  The stored binary64 entries of \(L\) are simply
interpreted as exact real constants in all subsequent validated calculations.
If the nominal factorization is inconvenient, another nonsingular positive-
diagonal scaling may be used; proof correctness does not depend on \(L\) being
an exact factor of \(A_c\).

Define the preconditioned matrix family
\[
  B(\pv)=L^{-1}A(\pv)L^{-T}.
\]
Because multiplication by a fixed nonsingular matrix is a congruence
transformation,
\[
  A(\pv)\succ0
  \quad\Longleftrightarrow\quad
  B(\pv)\succ0.
\]
Thus the proof problem can be moved to a family that is normally much closer to
the identity.

\subsection{Rigorous eigenvalue enclosure}

AutoIA first proves that the represented matrix family is structurally
symmetric.  Equal numerical interval endpoints alone do not prove that two
uncertain off-diagonal entries are the same expression; the symmetry proof must
retain the common dependency or prove their exact difference is zero.

For each row \(i\), let
\[
  B_{ii}(\pv)\in[d_i^-,d_i^+]
\]
be a rigorous enclosure of the diagonal entry and let
\[
  r_i
  \ge
  \sum_{j\ne i}\max_{\pv\in\Pin}|B_{ij}(\pv)|
\]
be a rigorous row-radius upper bound.  Gershgorin's theorem and symmetry imply
that every real eigenvalue \(\lambda\) of every admissible \(B(\pv)\) lies in
one of the intervals
\[
  [d_i^- - r_i,\; d_i^+ + r_i].
\]
Consequently the entire family satisfies
\[
  \lambda_L
  \le \lambda_k(B(\pv))\le
  \lambda_U,
\]
where
\[
  \lambda_L=\min_i(d_i^- - r_i),
  \qquad
  \lambda_U=\max_i(d_i^+ + r_i).
\]

AutoIA also forms
\[
  E(\pv)=B(\pv)-I
\]
and rigorously bounds a matrix norm by
\[
  \|E(\pv)\|_2\le\rho.
\]
Because \(B\) is symmetric, this independently gives
\[
  1-\rho\le\lambda_k(B(\pv))\le1+\rho.
\]
The two independent eigenvalue enclosures are intersected.  This is useful
because the near-identity norm theorem can be sharp locally while Gershgorin can
still prove positivity on boxes for which \(\rho\ge1\).

If the resulting lower bound satisfies
\[
  \boxed{\lambda_L>0},
\]
then
\[
  \boxed{A(\pv)\succ0\quad\text{for every }\pv\in\Pin.}
\]
If a sufficient bound is inconclusive, AutoIA may request physical-box
refinement.  Failure of one sufficient theorem is not interpreted as proof that
the matrix is non-SPD.

\subsection{Spectral log-determinant bounds}

For an SPD matrix, the determinant is positive.  Since
\[
  A=L B L^T,
\]
we have
\[
  \log\det A
  =2\sum_{i=1}^n\log L_{ii}+\log\det B.
\]
The eigenvalue enclosure immediately yields
\[
  n\log\lambda_L
  \le
  \log\det B
  \le
  n\log\lambda_U.
\]
Hence a broad rigorous interval is
\[
  \boxed{
  2\sum_i\log L_{ii}+n\log\lambda_L
  \le
  \log\det A(\pv)
  \le
  2\sum_i\log L_{ii}+n\log\lambda_U.
  }
\]
Every logarithm and reduction in this statement is evaluated with outward
rounding.

\subsection{Dependency-preserving matrix-log series}

The broad spectral interval can be tightened without discarding the shared
physical dependence.  Choose a positive stored binary64 scaling constant
\(c\) and write
\[
  B(\pv)=c\,[I+F(\pv)].
\]
AutoIA chooses \(c\) from the rigorous positive eigenvalue enclosure so that
\[
  \|F(\pv)\|_2\le\rho_s<1.
\]
For a symmetric matrix with spectral radius below one,
\[
  \log(I+F)
  =F-\frac{F^2}{2}+\frac{F^3}{3}-\cdots.
\]
Thus
\[
  \log\det B
  =n\log c
  +\operatorname{tr}\left(
  F-\frac{F^2}{2}+\cdots+(-1)^{r+1}\frac{F^r}{r}
  \right)+R_r,
\]
where \(r\) is the retained series order.  The retained matrix powers are
ordinary AutoIA expressions and therefore preserve dependence on \(\pv\).  The
omitted trace tail satisfies
\[
  |R_r|
  \le
  n\frac{\rho_s^{r+1}}{(r+1)(1-\rho_s)}.
\]
This gives a dependency-bearing log-determinant enclosure plus a rigorously
bounded scalar tail.  The final range is intersected with the independent
spectral log-determinant interval.

\subsection{Inverse quadratic forms without entrywise inversion}

A second ubiquitous quantity is
\[
  q(\pv)=b(\pv)^T A(\pv)^{-1}b(\pv),
\]
where \(b(\pv)\in\R^n\) may itself depend on the same uncertain inputs.  Define
\[
  z(\pv)=L^{-1}b(\pv).
\]
Then
\[
  q(\pv)=z(\pv)^T B(\pv)^{-1}z(\pv).
\]
The eigenvalue theorem gives the broad independent bounds
\[
  \boxed{
  \frac{\|z(\pv)\|_2^2}{\lambda_U}
  \le q(\pv)\le
  \frac{\|z(\pv)\|_2^2}{\lambda_L},
  }
\]
with rigorous interval bounds used for \(\|z(\pv)\|_2^2\).

The same scaling used for the matrix logarithm gives
\[
  B^{-1}=\frac{1}{c}(I+F)^{-1}
  =\frac{1}{c}\left[I-F+F^2-\cdots+(-F)^r\right]+\mathcal{R}_r.
\]
The inverse-series operator tail satisfies
\[
  \|\mathcal{R}_r\|_2
  \le
  \frac{1}{c}\frac{\rho_s^{r+1}}{1-\rho_s}.
\]
Therefore
\[
  \left|
  z^T\mathcal{R}_r z
  \right|
  \le
  \frac{\|z\|_2^2}{c}
  \frac{\rho_s^{r+1}}{1-\rho_s}.
\]
The retained expression
\[
  \frac{1}{c}z^T
  \left[I-F+F^2-\cdots+(-F)^r\right]z
\]
remains dependency-bearing in the original physical inputs.  AutoIA may also
attempt its independent general parameter-dependent linear-solve theorem and
intersect the resulting range if that theorem succeeds.  A failure of the
optional solve theorem cannot invalidate the already proved SPD series result.

\subsection{Worked two-by-two example}

Consider
\[
  A(x)=
  \begin{bmatrix}
  2+0.1x & 0.05x\\
  0.05x & 1.5-0.05x
  \end{bmatrix},
  \qquad x\in[-0.2,0.2],
\]
and
\[
  b(x)=
  \begin{bmatrix}
  1+0.2x\\ 0.5-0.1x
  \end{bmatrix}.
\]
All four uncertain matrix entries share the same scalar \(x\).  Replacing them
by four independent intervals would admit matrices that are not equal to
\(A(x)\) for any physical \(x\).  AutoIA instead retains the common \(x\), proves
the SPD family after constant preconditioning, and can certify
\[
  \log\det A(x),
  \qquad
  \det A(x),
  \qquad
  b(x)^T A(x)^{-1}b(x)
\]
with one structured theorem portfolio.  This is the general mechanism used by
V1; no covariance, Gaussian-process, or application-specific formula is built
into the matrix semantic.

%===============================================================================
\section{Validated implicit nonlinear solves}
%===============================================================================

Many engineering codes define an internal state implicitly rather than by an
explicit formula.  Let
\[
  F(y,p)=0
\]
be such a residual equation.  Here \(p\) denotes the physical uncertain inputs,
and \(y\) denotes the unknown state.  For a scalar root, \(y\in\mathbb{R}\); for
a square system, \(y\in\mathbb{R}^n\).

AutoIA deliberately separates the \emph{numerical proposal} from the
\emph{mathematical certificate}.  A recognized SciPy solver may supply a good
binary64 candidate \(y_0\), but SciPy convergence is never accepted as proof.
This design was chosen because the theorem that matters to the engineering user
is existence and enclosure of an exact real solution, not reproduction of every
floating-point iteration of Brent's method or MINPACK.

\subsection{Dependency-bearing implicit predictor}

Let
\[
  a(p)=(a_1(p),\ldots,a_r(p))^T
\]
collect the uncertain residual arguments actually reached by the residual.  At a
nominal representative \(a_0\) and nominal solution \(y_0\), define the proposal
Jacobians
\[
  J_y = \left.\frac{\partial F}{\partial y}\right|_{(y_0,a_0)},
  \qquad
  J_a = \left.\frac{\partial F}{\partial a}\right|_{(y_0,a_0)}.
\]
For a vector state, \(J_y\in\mathbb{R}^{n\times n}\) and
\(J_a\in\mathbb{R}^{n\times r}\).  If the nominal \(J_y\) is invertible, the
implicit-function sensitivity proposal is
\[
  S=-J_y^{-1}J_a.
\]
AutoIA then forms
\[
  \boxed{
  \widehat y(p)=y_0+S\bigl(a(p)-a_0\bigr).
  }
\]
The important feature is that \(a(p)\) remains an AutoIA value carrying its
original-input dependence.  The root is therefore not immediately replaced by
an independent scalar interval.

This choice is essential downstream.  For example, if
\[
  y(p)^2=p
\]
and a later line computes \(y^2-p\), a dependency-bearing predictor can preserve
near cancellation; an independent interval for \(y\) would generally destroy it.

\subsection{Bracketed scalar \texttt{root\_scalar}}

The current V1 semantic supports a finite constant bracket
\[
  [a,b],\qquad a<b,
\]
and a bracketed SciPy method.  AutoIA first proves a uniform endpoint sign
change.  If
\[
  F(a,p)\in A,
  \qquad
  F(b,p)\in B,
\]
then \(A\) and \(B\) must have strict opposite signs for every admissible
physical input.

AutoIA then proves uniqueness in the bracket by enclosing
\[
  F_y(y,p)=\frac{\partial F}{\partial y}(y,p)
\]
over the bracket.  If one bracket-wide enclosure is too loose, the root variable
is subdivided.  Every accepted subinterval must exclude zero and all accepted
subintervals must have the same derivative sign.  The endpoint sign change and
strict monotonicity then imply exactly one root in \([a,b]\) for every physical
input.

For local correction define
\[
  e=y-\widehat y(p)
\]
and choose the proposal preconditioner
\[
  C=J_y^{-1}.
\]
Consider
\[
  T(e,p)=e-CF(\widehat y(p)+e,p).
\]
AutoIA proves
\[
  |CF(\widehat y(p),p)|\le c
\]
and, for \(e\in[-d,d]\),
\[
  |1-CF_y(\widehat y(p)+e,p)|\le\rho.
\]
If
\[
  \rho<1,
  \qquad
  \frac{c}{1-\rho}\le d,
\]
then Banach's contraction theorem proves one correction in the tube and hence
\[
  \boxed{y(p)\in\widehat y(p)+[-d,d].}
\]
Because the bracket root was already proved unique, this local correction is the
same mathematical root requested by the bracketed call.

\paragraph{Example 1: square root as an implicit state.}
Let \(p\in[1,1.21]\) and
\[
  F(y,p)=y^2-p,
  \qquad y\in[0.5,2].
\]
The positive branch has exact solution \(y=\sqrt{p}\).  A nominal candidate near
\(p_0=1.105\) gives a sensitivity proposal
\[
  \frac{dy}{dp}\approx \frac{1}{2y_0}.
\]
AutoIA retains this first-order dependence and proves the remaining correction.

\paragraph{Example 2: Colebrook--White friction factor.}
The Darcy friction factor \(f\) is defined by
\[
  \frac{1}{\sqrt f}
  +2\log_{10}\!\left(
      \frac{\epsilon/D}{3.7}
      +\frac{2.51}{Re\sqrt f}
    \right)=0.
\]
Here \(Re\) is Reynolds number and \(\epsilon/D\) is relative roughness.
The canonical V1 example uses the user's ordinary \texttt{root\_scalar} call,
then AutoIA independently proves the bracket semantics and the dependency-bearing
correction.

\subsection{Square-system \texttt{fsolve}}

For
\[
  F(y,p)=0,
  \qquad y\in\mathbb{R}^n,
\]
the V1 \texttt{fsolve} theorem is intentionally local because \texttt{fsolve}
itself is an initial-guess local solver.  With a finite proposal
\(C\approx J_y^{-1}\), AutoIA proves
\[
  \|CF(\widehat y(p),p)\|_\infty\le c
\]
and, over \(\|e\|_\infty\le d\), a componentwise defect matrix
\[
  B\ge |I-CJ_y(\widehat y(p)+e,p)|.
\]
Let
\[
  \rho=\|B\|_\infty.
\]
If
\[
  \rho<1,
  \qquad
  \frac{c}{1-\rho}\le d,
\]
then a unique solution exists in that correction tube for every physical input.
This theorem does not claim that no remote solution exists elsewhere.

\paragraph{Example 3: coupled two-state equilibrium.}
For uncertain \(a,b\), consider
\[
  F_1(y,a,b)=y_1+0.2y_2-a,
\]
\[
  F_2(y,a,b)=0.1y_1+y_2+0.05y_1y_2-b.
\]
The nominal solve proposes \(y_0\), the implicit Jacobian constructs a predictor,
and the contraction proof validates a local unique branch.

\subsection{Derivative trust boundary}

Forward-mode differentiation is used to construct nominal and proof expressions.
A nominal derivative is only proposal information.  Proof-critical derivative
ranges are themselves validated through the supported scalar language and trusted
range machinery.  A derivative through an unsupported operation causes an
explicit failure rather than a floating-point finite-difference substitute.


%===============================================================================
\section{Generator-aware semantic linear solve}
%===============================================================================

Suppose the solve inputs contain shared auxiliary generators:
\[
  A(\qv,\epsv)
  \in
  \widehat A(\qv)+\sum_{j=1}^s A_j\eps_j+R_A,
\]
\[
  b(\qv,\epsv)
  \in
  \widehat b(\qv)+\sum_{j=1}^s b_j\eps_j+R_b.
\]
AutoIA proposes a solution of the form
\[
  \boxed{
  X(\qv,\epsv)
  =P(\qv)+\sum_{j=1}^s u_j\eps_j.
  }
\]
For generator \(\eps_j\), a useful sensitivity proposal is obtained from a
reference matrix \(A_{\mathrm{ref}}\) by
\[
  A_{\mathrm{ref}}u_j
  =b_j-A_jx_0,
\]
where \(x_0=P(0)\).  Again, this is only a proposal.

The proof then evaluates the actual residual
\[
  r(\qv,\epsv)
  =b(\qv,\epsv)-A(\qv,\epsv)X(\qv,\epsv)
\]
over the joint domain
\[
  \qv\in\Qdom,
  \qquad
  \epsv\in[-1,1]^s,
\]
and uses the same fixed-point theorem.  Incoming generator identifiers are
preserved in the explicit predictor when justified.  The remaining componentwise
correction proven by the contraction theorem is assigned fresh independent
auxiliary generators because the theorem proves only a box for the residual
error, not a stronger coupled geometry.

%===============================================================================
\section{Path-aware execution and uncertain control flow}
%===============================================================================

Consider a scalar predicate such as
\[
  \texttt{lhs < rhs}.
\]
AutoIA forms the validated difference
\[
  d=\texttt{lhs}-\texttt{rhs}
\]
and computes a rigorous range
\[
  d\in[\underline d,\overline d].
\]
For the strict less-than relation:
\begin{itemize}[leftmargin=*]
\item if \(\overline d<0\), the predicate is proved true throughout the current
box;
\item if \(\underline d\ge0\), it is proved false throughout the current box;
\item otherwise the predicate is unresolved.
\end{itemize}
Analogous exact-real rules are used for \(\le,>,\ge,==,!=\) and scalar truth
checking.

If the predicate is unresolved, AutoIA may split the physical input box and
re-execute.  For an exact affine one-variable predicate
\[
  d(p_i)=a+bp_i,
  \qquad b\ne0,
\]
the exact root
\[
  p_i^\star=-a/b
\]
can be used as a direct partition point when it is rigorously recognized.

For a general nonlinear or multivariate decision surface, finitely many Cartesian
boxes may not align with the exact path boundary.  If subdivision cannot resolve
the predicate under the allowed budget, AutoIA conservatively enumerates both
feasible Boolean outcomes and certifies the continuation of each path separately.
If the two branch enclosures are \(Y_T\) and \(Y_F\), then
\[
  \mathcal{P}(\Pin)\subseteq Y_T\cup Y_F.
\]
The branch formula may be evaluated on a superset of its true feasible domain;
this can make the result wider or cause a boundary singularity to fail
certification, but it cannot omit a true output.

Path-local continuation is essential.  AutoIA does not immediately hull two
branch values and then continue the rest of the program.  Instead every feasible
path carries its own validated values through downstream arithmetic, elementary
functions, and semantic solves.  The union is formed only at the final piecewise
result.

%===============================================================================
\section{Exact closed-endpoint comparison metadata}
%===============================================================================

Outward rounding is required for rigorous numerical enclosures, but it can make
an expression that is mathematically exact at a closed endpoint appear to cross
a comparison threshold by a few ulps.  For example, if
\[
  c\in[0.95,1]
\]
then the physical statement \(c>1\) is false everywhere, even if a rounded
polynomial range of \(c-1\) is
\[
  [-0.05,7\times10^{-18}].
\]

To preserve exact Python comparison semantics in such cases, AutoIA carries a
\emph{proof-only exact affine shadow} whenever an expression remains exactly
affine in the original physical inputs.  Write
\[
  s(p)=\beta_0+\sum_{i=1}^m\beta_i p_i,
\]
where each \(\beta_i\) is stored as an exact rational number equal to the exact
real value of the corresponding binary64 constant.  Over a box, the exact range
of this affine function is obtained by choosing each interval endpoint according
to the sign of \(\beta_i\).  This metadata is used only to decide comparisons;
it never replaces the trusted numerical enclosure used for arithmetic.

The shadow survives addition, subtraction, negation, multiplication by an exact
constant, and division by an exact nonzero constant.  Nonlinear operations
discard it.  This design was chosen because it solves endpoint semantics without
weakening outward rounding or introducing symbolic algebra into general
nonlinear arithmetic.

\paragraph{Example 1.}
For \(x\in[0.95,1]\), the branch
\begin{verbatim}
if x > 1.0:
    raise ValueError(...)
\end{verbatim}
is proved infeasible exactly.

\paragraph{Example 2.}
For the same box, the predicates \(x<1\) and \(x==1\) define two feasible path
pieces.  AutoIA retains the equality boundary rather than silently assigning it
to the wrong strict branch.

%===============================================================================
\section{Proof-driven physical-domain refinement}
%===============================================================================

A broad enclosure may fail to prove a domain condition even when the original
engineering model is valid throughout the physical box.  AutoIA therefore treats
several domain obligations as requests for refinement rather than immediate
terminal failures.

For a reciprocal
\[
  z=\frac{1}{d(p)},
\]
the proof obligation is
\[
  0\notin d(\Pin).
\]
If the current enclosure of \(d\) contains zero but does not prove an actual
zero, the controller may subdivide the physical box, re-execute the model, and
prove the denominator nonzero on each child.  The same architectural pattern is
used for supported elementary-domain conditions such as positivity of a logarithm
argument.

This is sound because no parent proof is accepted.  Each retained child is
independently certified and the child boxes cover the parent domain.

\paragraph{Example 1: removable over-width in a denominator.}
A denominator may be positive on the true domain while dependency loss in a broad
intermediate enclosure crosses zero.  Two physical subboxes can recover the
positive proof without changing the engineering formula.

\paragraph{Example 2: shell-and-tube heat exchanger algebra.}
A public effectiveness formula contains nested reciprocals.  On a broad
\((NTU,C_r)\) box, a whole-box denominator enclosure can cross zero even though
finer physical boxes certify.  The V1 controller now requests reciprocal-driven
subdivision before failing closed.

This mechanism does not perform symbolic cancellation of removable singularities.
If a branch formula genuinely approaches a \(0/0\) form on the closure of a
strict path, generic interval division can still fail.  That is a separate
mathematical issue.

%===============================================================================
\section{Fixed-data regular-grid multilinear interpolation}
%===============================================================================

The V1 semantic for \texttt{scipy.interpolate.RegularGridInterpolator} targets
fixed numeric rectilinear grids and fixed numeric table values with
\texttt{method="linear"}.  Uncertainty is allowed in the query coordinates.
This subset was chosen because tabulated maps are ubiquitous in engineering and
piecewise multilinear interpolation has a transparent exact algebraic meaning.

Let there be \(d\) grid dimensions.  On dimension \(i\), let the strictly
monotone grid coordinates be
\[
  g_{i,0},g_{i,1},\ldots,g_{i,n_i-1}.
\]
After internal normalization to ascending order, suppose a query coordinate
\(x_i\) lies in cell
\[
  g_{i,j_i}\le x_i\le g_{i,j_i+1}.
\]
Define the normalized local coordinate
\[
  t_i=\frac{x_i-g_{i,j_i}}{g_{i,j_i+1}-g_{i,j_i}}.
\]
For a scalar table, the multilinear interpolant in that cell is
\[
  \boxed{
  I(x)=
  \sum_{b\in\{0,1\}^d}
  v_{j_1+b_1,\ldots,j_d+b_d}
  \prod_{i=1}^d
  \bigl[b_i t_i+(1-b_i)(1-t_i)\bigr].
  }
\]
The vector \(b=(b_1,\ldots,b_d)\) selects one of the \(2^d\) cell corners.
Because the tabulated values are fixed constants, this is ordinary AutoIA
arithmetic in the uncertain query coordinates and preserves their dependence.

If a query box crosses a grid boundary, the boundary comparison becomes an
ordinary AutoIA path decision.  Every feasible cell formula is certified on its
path and the final piecewise hull covers all cells reached by the physical input
box.  No global min/max lookup is substituted for the interpolant.

\paragraph{Example 1: bilinear table.}
For a two-dimensional cell with corner values \(v_{00},v_{10},v_{01},v_{11}\),
\[
 I(x_1,x_2)=(1-t_1)(1-t_2)v_{00}
 +t_1(1-t_2)v_{10}
 +(1-t_1)t_2v_{01}
 +t_1t_2v_{11}.
\]
If the uncertain box remains inside the cell, this one expression is propagated
through downstream code.

\paragraph{Example 2: aerodynamic lookup crossing cells.}
Suppose Mach number and angle of attack are uncertain and both cross table
breakpoints.  AutoIA splits/enumerates the corresponding cell paths, evaluates
the proper bilinear formula on each, and hulls only at the final path result.

\paragraph{Example 3: out-of-domain semantics.}
With \texttt{bounds\_error=True}, a physical box containing an out-of-grid query
means the user model is not defined everywhere and certification fails.  With
\texttt{bounds\_error=False} and a finite fixed \texttt{fill\_value}, the outside
path returns that constant.  With \texttt{fill\_value=None}, the first or last
cell formula supplies the same linear extrapolation semantics as SciPy.

Nonlinear interpolation methods, uncertain grid coordinates, and uncertain table
values are deliberately not part of this V1 semantic.


%===============================================================================
\section{Adaptive subdivision controller}
%===============================================================================

Subdivision is a sharpness/resource strategy, not a correctness prerequisite.
Every retained child box is independently certified.  If a controller chooses a
poor direction, the result may be slower or wider but does not become unsound.

The practical split score introduced in version 0.19.0.dev1 uses the already
certified parent polynomial.  For candidate coordinate \(q_d\), a half-box child
has a shifted center
\[
  s_d=\pm h_d/2
\]
and new coordinate \(r_d\) satisfying
\[
  q_d=r_d+s_d,
  \qquad
  r_d\in[-h_d/2,h_d/2].
\]
The binomial identity
\[
  q_d^{\alpha_d}
  =\sum_{b=0}^{\alpha_d}
    {\alpha_d\choose b}s_d^{\alpha_d-b}r_d^b
\]
permits a cheap predicted re-centering without executing the engineering model.
The score estimates derivative-sign resolution, reduced-Bernstein burden, and
output-width gain.  Only the selected child pair is actually certified.  The
controller then checks realized structural/width improvement before accepting
the split.

The old expensive one-step lookahead remains available as a research oracle.
Because controller decisions are outside the trusted proof kernel, agreement with
that oracle is an efficiency/sharpness validation issue rather than a soundness
condition.

%===============================================================================
\section{Automatic configuration as a portfolio of independent proofs}
%===============================================================================

The novice API
\[
  \texttt{autoia(model, box)}
\]
starts from a release baseline consisting of the polynomial representation at
order four.  The integer \(k=4\) here is the maximum retained total polynomial
degree.  The controller may make a bounded number of optional retries when the
baseline exposes one of two generic structural mechanisms:
\begin{enumerate}[leftmargin=*]
\item the rigorous remainder dominates the retained physical-polynomial
variation, suggesting that a modestly higher order may help; or
\item approximation-error ancestry is reused/recombined downstream, suggesting
that shared generators may preserve cancellation.
\end{enumerate}
These triggers are heuristics for cost and sharpness only.

Suppose successful independent runs produce output intervals
\[
  Y^{(1)},Y^{(2)},\ldots,Y^{(k)}
\]
for the same scalar output.  Because every interval contains the exact output,
the intersection
\[
  \boxed{Y=\bigcap_{j=1}^kY^{(j)}}
\]
also contains the exact output.  Componentwise, the reported bounds are
\[
  \underline Y=\max_j\underline Y^{(j)},
  \qquad
  \overline Y=\min_j\overline Y^{(j)}.
\]
Thus the selector is not trusted proof arithmetic.  A poor selector can waste
work or miss an optional tightening; it cannot make the intersection unsound.

\paragraph{Example.}
If a polynomial/order-4 run proves \([-2,5]\) and a generator/order-4 retry proves
\([-1,6]\), the intersection is \([-1,5]\), which is rigorous even though
neither run dominates the other completely.

%===============================================================================
\section{Application-theorem interface}
%===============================================================================

AutoIA's generic execution engine should not absorb application-specific
mathematics.  Instead, an application expert may state a theorem externally and
use a small rigorous context to evaluate its numerical consequences.

Let a theorem prove that a scalar target \(f\) is nondecreasing in a certified
intermediate \(z\) on the current domain.  If AutoIA encloses
\[
  z(p)\in[\underline z,\overline z],
\]
then the theorem may conclude
\[
  f(p)\le g(\overline z)
\]
for the specified monotone outer function \(g\).  The application layer records
which assumptions were checked and returns a typed rigorous upper or lower bound.

Independent application theorems can also be combined.  If
\[
  f(p)\le U_1
  \qquad\text{and}\qquad
  f(p)\le U_2,
\]
then
\[
  f(p)\le \min(U_1,U_2).
\]
Similarly, independent lower bounds combine by a maximum, and two full certified
intervals for the same quantity combine by intersection.

\paragraph{Example 1: monotonicity.}
If \(f(x)=e^x\) and \(x\in[\ell,u]\), monotonicity gives
\[
  f(x)\le e^u.
\]
The theorem is supplied by the application expert; the exponential evaluation is
performed rigorously by AutoIA.

\paragraph{Example 2: Lipschitz theorem.}
If the expert has proved
\[
  |f(x)-f(c)|\le L\|x-c\|_2
\]
on the application domain, AutoIA can rigorously evaluate \(f(c)\), the box
radius, and the resulting upper/lower bound without embedding the theorem in the
core engine.

This separation is the permanent architectural boundary used by application
packages such as downstream GP application: application mathematics may depend on AutoIA, while
AutoIA itself remains application-independent.


%===============================================================================
\section{Prepared sessions and repeated-query certification}
%===============================================================================

Version 1.0.2 retains the V1 separation of static program preparation from repeated box
certification.  Let one fixed engineering program \(\mathcal{P}\) and one fixed
AutoIA configuration be used for a sequence of boxes
\[
  \Pin^{(1)},\Pin^{(2)},\ldots.
\]
A prepared session reuses the local semantic-dispatch clone and one audited
trusted backend.  More importantly, it stores previously certified boxes and
results in a bounded certificate cache.

\subsection{Exact repeated-box reuse}

If a new query box has exactly the same stored binary64 endpoints as a cached
box, the previous certificate already proves the same statement.  No engineering
model execution is required.

\subsection{Restriction of a parent certificate to an interior subbox}

Suppose a parent box \(D\) has already been certified and a new query box \(Q\)
satisfies
\[
  Q\subseteq D.
\]
Set-theoretically, every certificate valid on \(D\) is automatically valid on
\(Q\).  For a polynomial certificate, AutoIA can do better than returning the
same parent output interval: it re-centers and re-ranges the stored polynomial on
\(Q\).

Let the old centered coordinate be
\[
  \qv=\pv-\mathbf{c}_{\mathrm{old}},
\]
and the new centered coordinate be
\[
  \mathbf{r}=\pv-\mathbf{c}_{\mathrm{new}}.
\]
Then exactly
\[
  \qv=\mathbf{r}+\mathbf{s},
\]
where
\[
  \mathbf{s}=\mathbf{c}_{\mathrm{new}}-\mathbf{c}_{\mathrm{old}}.
\]
For a monomial
\[
  \qv^\alpha=\prod_{i=1}^m(r_i+s_i)^{\alpha_i}.
\]
Using the multivariate binomial expansion,
\[
  \qv^\alpha
  =
  \sum_{\beta\le\alpha}
  \left[
  \prod_{i=1}^m
  \binom{\alpha_i}{\beta_i}
  s_i^{\alpha_i-\beta_i}
  \right]
  \mathbf{r}^\beta.
\]
Therefore if
\[
  P(\qv)=\sum_\alpha c_\alpha\qv^\alpha,
\]
the exact re-centered coefficient of \(\mathbf{r}^\beta\) is
\[
  c'_\beta
  =
  \sum_{\alpha\ge\beta}
  c_\alpha
  \prod_{i=1}^m
  \binom{\alpha_i}{\beta_i}
  s_i^{\alpha_i-\beta_i}.
\]

The exact shift \(s_i\) is not assumed to equal the ordinary floating-point
subtraction of the two stored centers.  The trusted backend forms an interval
\(S_i\) containing the exact real difference.  Every transformed coefficient is
computed as a rigorous interval using \(S_i\), while an ordinary binary64 point
coefficient is stored as a proposal.  The difference between the exact coefficient
interval and the stored proposal is rigorously ranged over the new centered box
and accumulated into an extra remainder.

Hence the re-centered value remains a valid certificate on the new domain.
Shared-generator coefficients are copied unchanged because the coordinate shift
acts on the physical polynomial \(P\), not on the auxiliary variables \(\epsv\).

\subsection{Containment check for the actual centered query domain}

The user query \(Q\) is not quite the domain used by the polynomial executor.
The executor constructs an outward symmetric centered envelope around each new
binary64 center.  Let that actual represented domain be \(\widetilde Q\).
Certificate restriction is used only when AutoIA proves
\[
  \widetilde Q\subseteq D.
\]
This prevents the known one-ulp symmetric-centering enlargement from causing a
parent certificate to be applied outside the domain on which it was proved.  If
this containment check fails, AutoIA performs a fresh certification instead.

\subsection{Many-box execution and multicore CPU proof}

For a collection of \(N\) independent physical boxes, prepared execution supports
serial, threaded, and process-parallel modes.  With
\texttt{workers="auto"} and \texttt{parallel="auto"}, AutoIA resolves an
available CPU worker count and prefers process parallelism when the model can be
serialized.  Optional \texttt{cloudpickle}, when installed, permits many local
functions and closures to be serialized for spawn-based workers; it is execution
infrastructure and is not part of the proof.

In process mode each worker constructs one prepared session and one independently
audited trusted backend, then handles multiple boxes.  The mathematical statement
proved for box \(B_k\) is independent of which worker evaluated it.  Parallel
scheduling therefore changes wall time but not the inclusion theorem.  The public
many-box call records non-proof execution metadata such as the resolved worker
count, serial/thread/process mode, chunk size, elapsed time, and model serializer.

A streaming interface keeps only a bounded number of chunks in flight, so the
user does not need to store all \(N\) boxes or all \(N\) full result objects in
memory simultaneously.  This is execution infrastructure for large repeated-query
workloads; it does not decide which boxes an optimization or reliability algorithm
should generate.

\subsection{MPS and CUDA proposal acceleration with CPU verification}

AutoIA V1 also exposes optional Torch-backed proposal helpers for Apple Metal
through the PyTorch MPS device and for NVIDIA hardware through CUDA.  These
accelerator calculations are explicitly outside the trusted proof base.  A GPU
may evaluate many center points, ranking scores, nominal preconditioners, or other
proposal quantities, but the resulting object carries
\texttt{machine\_rigorous=False}.  Any final AutoIA certificate is recomputed or
verified independently by the trusted CPU binary64 backend.

The V1 trust boundary is therefore
\[
  \boxed{
  \text{optional MPS/CUDA proposal or assist}
  \quad\longrightarrow\quad
  \text{trusted multicore CPU verification}.
  }
\]
This allows applications to exploit available accelerators without making the
correctness of an enclosure depend on GPU rounding modes, tensor-library kernel
choices, or hardware-specific reduced precision.  A future fully trusted GPU
backend would require a separate directed-rounding implementation and audit.

%===============================================================================
\section{Program-level soundness theorem}
%===============================================================================

We can summarize the architecture as a compositional theorem.

\paragraph{Assumptions.}
Assume:
\begin{enumerate}[leftmargin=*]
\item the trusted binary64 backend's audited machine assumptions hold;
\item every proof-critical interval operation returns an enclosure of its exact
real counterpart;
\item every supported validated scalar/semantic operation satisfies its enclosure
contract;
\item every representation conversion is inclusion-preserving;
\item every retained path/subdivision child is certified over a domain whose
union covers the relevant user input domain;
\item every semantic linear or nonlinear solve accepted by the system satisfies
its stated residual/contraction and, where applicable, bracket-uniqueness theorem;
\item every accepted SPD matrix-family operation satisfies its structural symmetry,
positive-eigenvalue, matrix-series, and tail-bound proof obligations;
\item every accepted fixed-data interpolation path uses the correct multilinear
cell formula on a domain covering that path; and
\item the program remains within the documented supported-language semantics.
\end{enumerate}

\paragraph{Conclusion.}
For every user input \(\pv\in\Pin\), the exact real output of the supported
program semantics belongs to the finite union of the validated output sets
returned by AutoIA.  Consequently the reported global lower/upper arrays contain
every exact scalar output.

\paragraph{Proof sketch.}
The proof is by induction through execution.  Initial uncertain inputs are
represented by enclosures containing every physical value in the input box.
Every supported operation maps valid input enclosures to an output enclosure
containing the exact operation on all admissible inputs.  Semantic solves are
validated independently by their residual/contraction theorem.  If execution
branches, the path engine retains every feasible outcome or subdivides the
physical domain into certified children whose union covers the original domain.
Thus the invariant holds on every active path.  Final range methods enclose each
representation's concretization, and the finite union/hull of path results
contains the exact output set.

Prepared-session restriction does not alter the theorem: a parent certificate is
reused only on a subdomain proved to lie inside the parent's certified domain,
and the coordinate transform itself is rigorously enclosed.

%===============================================================================
\section{What ordinary floating-point calculations are allowed to do}
%===============================================================================

A recurring AutoIA design principle is that an approximation may be poor without
making the final result unsound.  Ordinary floating point is therefore acceptable
for:
\begin{itemize}[leftmargin=*]
\item polynomial coefficient proposals;
\item nominal/reference linear solves;
\item approximate inverses and preconditioners;
\item nominal Cholesky factors and constant SPD preconditioners;
\item midpoint representatives of interval matrix families;
\item branch split scores and controller heuristics;
\item candidate Taylor centers;
\item integer argument-reduction proposals for trigonometric functions;
\item cache-choice heuristics such as choosing the smallest containing parent
certificate.
\end{itemize}

Such calculations may affect sharpness, performance, or whether a proof attempt
succeeds.  They must not be used as the final reason an inclusion is accepted.
The proof always returns to rigorous residual, interval, range, or domain checks.

%===============================================================================
\section{V1 heterogeneous execution policy}

Execution scheduling is outside the mathematical proof.  For the preferred finite many-box interface, the default \texttt{workers="auto"}, \texttt{parallel="auto"} policy stays serial for fewer than four chunks and otherwise uses at most one CPU worker per chunk, capped by the available CPU count.  Explicit user scheduling overrides the heuristic.  The resulting certificates are unchanged in meaning because every box is independently verified by the trusted CPU binary64 backend.

Optional MPS and CUDA helpers are proposal/assist devices only.  Version 1 deliberately resolves \texttt{device="auto"} to CPU rather than guessing GPU profitability from batch size.  The application must explicitly request \texttt{mps} or \texttt{cuda}.  This is a performance policy and does not alter the trust boundary: GPU results are never machine-rigorous certificates.

\section{Important V1 limitations and explicit non-guarantees}
%===============================================================================

AutoIA is a validated-execution system for a deliberately bounded Python/NumPy/
SciPy language surface.  The following are current non-claims.

\begin{enumerate}[leftmargin=*]
\item \textbf{Exact hulls are not promised.}  Every accepted interval is an
inclusion certificate, not necessarily the tightest possible interval.
\item \textbf{Arbitrary Python is not supported.}  General side effects,
uncertainty-dependent indexing, arbitrary Boolean masks, unrestricted
uncertainty-dependent loop termination, and opaque compiled-extension internals
are outside the frozen V1 language unless a separate semantic primitive exists.
\item \textbf{Local imports are a dispatch boundary.}  The source-level shadow
namespace prepares visible globals and reached source-backed helpers; imports
performed dynamically inside a function body are not generally rewritten.
\item \textbf{Shared generators are affine auxiliary variables.}  Products such
as \(q_i\varepsilon_j\) and \(\varepsilon_i\varepsilon_j\) are bounded back into
new rigorous error rather than retained as a richer nonlinear auxiliary basis.
\item \textbf{Generator compression is not yet general.}  Long calculations can
accumulate auxiliary generator identifiers.
\item \textbf{Implicit solves are conditional semantics.}  Bracketed
\texttt{root\_scalar} requires a uniform bracket/uniqueness proof;
\texttt{fsolve} proves only a local unique branch around the validated tube.
Bifurcations, singular Jacobians, or merged branches can fail closed.
\item \textbf{Interpolation support is intentionally narrow.}  V1 certifies
fixed-data \texttt{RegularGrid\allowbreak Interpolator} queries with \texttt{method="linear"}.  Nonlinear
methods and uncertain table/grid data are not claimed.
\item \textbf{General eigensystems are not yet certified.}
\texttt{scipy.linalg.eigh} with AutoIA-dependent matrices is recognized and fails
explicitly with \texttt{UnsupportedOperationError}; no eigenvalue or eigenvector
certificate is produced.
\item \textbf{General ODE/DAE and optimizer semantics are not yet certified.}
Calls such as \texttt{solve\_ivp} or arbitrary constrained optimization remain
outside the V1 semantic set unless a future version adds a theorem-backed
primitive.
\item \textbf{Removable singularities are not symbolically regularized.}  A
branch formula whose closure contains a true \(0/0\) form may fail even if the
piecewise mathematical function has a finite limit.
\item \textbf{GPU proof execution is not part of V1.}  V1 exposes optional
PyTorch MPS and CUDA proposal/assist helpers, but these results are explicitly
non-rigorous.  The current trusted proof kernel, including parallel many-box
verification, is CPU binary64.
\item \textbf{Piecewise cached restriction and giant many-box kernels remain
limited.}  Prepared sessions provide exact and contained-box reuse, but not a
fully general symbolic restriction engine for arbitrary piecewise branch
geometry or one dense GPU-sized batch proof kernel.
\end{enumerate}

A recognized unsupported operation must fail closed.  The desired V1 behavior is
therefore not ``every engineering program certifies,'' but rather ``supported
semantics certify rigorously, and unsupported or unproved semantics are made
explicit without silently returning an uncertified result.''

%===============================================================================
%===============================================================================
\section{A compact end-to-end mathematical example}
%===============================================================================

Consider a two-parameter model
\[
  p_1\in[0.9,1.1],
  \qquad
  p_2\in[-0.1,0.1],
\]
with
\[
  A(\pv)=
  \begin{bmatrix}
  2+p_1^2 & 0.2\\
  -0.1p_2 & 1.5+0.05p_1
  \end{bmatrix},
\]
\[
  b(\pv)=
  \begin{bmatrix}
  1+0.1p_2\\
  0.5+0.05p_1^2
  \end{bmatrix},
\]
and let \(u(\pv)\) solve
\[
  A(\pv)u(\pv)=b(\pv).
\]
Suppose the program returns
\[
  y_1=u_1,
  \qquad
  y_2=e^{u_2},
  \qquad
  y_3=u_1-u_1.
\]

AutoIA first constructs rigorous centered coordinates for \(p_1,p_2\).  Polynomial
arithmetic retains the quadratic occurrences of \(p_1\) up to the chosen order.
The semantic solve constructs a polynomial (and, in generator mode, auxiliary
sensitivity) predictor, then proves a contraction bound on the complete matrix
family.  The exponential is composed by the validated Taylor machinery.  Shared
generator provenance can preserve exact downstream cancellation in \(y_3\), so
\(y_3\) can certify as zero even when intermediate approximations carry error.

If a prepared session first certifies the master box and later receives a smaller
interior box, the stored output polynomial certificates may be re-centered using
the exact binomial coordinate transform with rigorous coefficient-error capture.
The model and the semantic solve need not run again if the new centered domain is
proved to lie inside the parent domain.

%===============================================================================
\section{Additional worked examples connecting the concepts}
%===============================================================================

\subsection{Example A: polynomial truncation versus interval dependence loss}

Let
\[
  f(x)=x-x+x^2,
  \qquad x\in[-1,1].
\]
The exact function is \(x^2\in[0,1]\).  A dependency-blind interval calculation
of the first two terms could form
\[
  [-1,1]-[-1,1]=[-2,2]
\]
before adding \([0,1]\).  A physical-input polynomial representation keeps the
same symbolic coordinate \(q\) in both occurrences and cancels the linear terms
before final ranging.

\subsection{Example B: generator cancellation after approximation}

Suppose an elementary approximation produces
\[
  y=P(q)+g\varepsilon+R,
  \qquad \varepsilon\in[-1,1].
\]
If downstream code returns \(y-y\), shared use of the same generator identifier
cancels \(g\varepsilon\) exactly.  Replacing both occurrences by independent
interval remainders would not preserve that fact.

\subsection{Example C: path-aware absolute-value-like code}

Consider
\begin{verbatim}
if x >= 0:
    y = x
else:
    y = -x
return y + x*x
\end{verbatim}
for \(x\in[-1,1]\).  AutoIA retains two path-local expressions and ranges them
separately.  It does not first hull \(y\) to \([0,1]\) and then forget which sign
of \(x\) produced it.

\subsection{Example D: semantic linear solve}

For
\[
  A(k)=\begin{bmatrix}2k&-k\\-k&k\end{bmatrix},
  \qquad
  b(F)=\begin{bmatrix}F\\0\end{bmatrix},
\]
with uncertain stiffness \(k\) and load \(F\), NumPy's nominal
\texttt{linalg.solve} is only a proposal.  AutoIA verifies the matrix family by a
preconditioned residual/contraction theorem and returns a dependency-bearing
solution enclosure.

\subsection{Example E: implicit solve followed by dependent arithmetic}

Let \(y(p)\) be the positive solution of \(y^2=p\).  If downstream code computes
\[
  r=y^2-p,
\]
the exact value is zero.  A root interval alone loses the relation between \(y\)
and \(p\); the implicit predictor \(\widehat y(p)+[-d,d]\) preserves enough shared
dependence to make the residual enclosure far smaller.

\subsection{Example F: interpolation across a grid boundary}

Let a one-dimensional table have nodes \(g_0<g_1<g_2\), and suppose
\(x\in[g_0+\delta,g_2-\delta]\) crosses \(g_1\).  The interpolant is one affine
formula on \([g_0,g_1]\) and another on \([g_1,g_2]\).  AutoIA treats \(x\le g_1\)
as a path boundary, certifies each affine continuation, and hulls the two final
outputs.  This is exact with respect to the piecewise-linear interpolation
semantics apart from outward arithmetic rounding.

\subsection{Example G: independent theorem portfolio}

Suppose a generic execution proves \(f\in[-3,7]\), a monotonicity theorem proves
\(f\le5\), and a Lipschitz theorem proves \(f\ge-2\).  Combining independent
certificates gives
\[
  f\in[-2,5].
\]
No theorem needs to know how the others were obtained.

\subsection{Example H: correct fail-closed behavior for modal analysis}

For a structural problem
\[
  K(p)\phi=\lambda M(p)\phi,
\]
V1 recognizes \texttt{scipy.linalg.eigh} but does not yet possess the spectral
separation/eigenspace theorem needed for a general certificate.  The correct
V1 result is therefore an explicit unsupported-operation error, not a nominal
frequency with a misleading ``rigorous'' label.


\subsection{Example I: why an SPD matrix family is not an entrywise interval matrix}

Let
\[
  A(t)=\begin{bmatrix}1+t&t\\t&1+t\end{bmatrix},
  \qquad t\in[0,0.4].
\]
The same scalar \(t\) appears in all four uncertain locations.  Entrywise
intervalization would admit, for example, a matrix with one off-diagonal entry
zero and the other equal to \(0.4\), although no physical value of \(t\) creates
that matrix.  The structured matrix-family representation excludes such fictitious
combinations and permits preconditioned SPD reasoning over the actual one-parameter
family.

\subsection{Example J: inverse quadratic form by a dependency-bearing series}

Suppose a preconditioned two-by-two family has
\[
  B(s)=c[I+F(s)],\qquad \|F(s)\|_2\le0.2.
\]
Retaining terms through order two gives
\[
  B(s)^{-1}\approx\frac1c[I-F(s)+F(s)^2].
\]
The omitted operator norm is no larger than
\[
  \frac1c\frac{0.2^3}{1-0.2}=\frac{0.01}{c}.
\]
For \(z(s)=L^{-1}b(s)\), the retained quadratic expression
\[
  \frac1c z(s)^T[I-F(s)+F(s)^2]z(s)
\]
still contains the original uncertain variable \(s\); only the remaining tail is
collapsed to a scalar rigorous error interval.

\subsection{Example K: parallel boxes do not create a parallel-proof theorem}

Let three boxes \(B_1,B_2,B_3\) be certified on three CPU worker processes.  If
worker \(k\) proves
\[
  f(B_k)\subseteq Y_k,
\]
then the collection \(\{Y_1,Y_2,Y_3\}\) is valid regardless of execution order.
No theorem combines floating-point states across workers.  Parallelism is only a
scheduler around three independent proof calls.

\subsection{Example L: GPU proposal versus CPU proof}

Suppose an MPS kernel evaluates a box center and returns the float32 proposal
\[
  \widehat f(c)=1.2374.
\]
This number is useful for ranking or initialization but proves nothing.  AutoIA
then independently certifies the full physical box on CPU and obtains
\[
  f(B)\subseteq[1.18,1.31].
\]
Only the latter interval is a certificate.  A wrong GPU proposal may reduce
performance, but it cannot invalidate the CPU inclusion proof.

%===============================================================================
\section{Summary of the current mathematical architecture}
%===============================================================================

The current architecture can be summarized by the chain
\[
\boxed{
\begin{array}{c}
\text{user interval box}\\
\downarrow\\
\text{rigorous centered physical coordinates }\qv\\
\downarrow\\
P(\qv)+R\quad\text{or}\quad P(\qv)+G\epsv+R\\
\downarrow\\
\text{validated generic operations + semantic primitives}\\
\downarrow\\
\text{path-aware piecewise execution when required}\\
\downarrow\\
\text{termwise / monotonicity / Bernstein range certification}\\
\downarrow\\
\text{rigorous binary64 output intervals}
\end{array}}
\]
with a prepared-session layer that can reuse the static dispatch architecture and,
more importantly, restrict already-proved certificates to eligible interior
query boxes.

The central theoretical thesis is that sharp automatic interval analysis for
engineering code is best treated as a structured validated-execution problem:
retain original-input dependence explicitly, introduce auxiliary variables only
when their provenance is proved, treat important numerical operations according
to their mathematical semantics, and confine correctness-critical floating-point
work to a small audited enclosure kernel.  Performance strategies are then free
to be aggressive because any proposed approximation must still pass an independent
rigorous certificate before it is accepted.


\section{Transparent fitted exact-GP posterior semantics in Version 1.0.2}

Let $x\in\mathbb{R}^d$ denote an uncertain query vector with $d$ input
coordinates.  For a supported fitted scalar exact Gaussian process, AutoIA
freezes the numerical fitted state and reconstructs the posterior using validated
arithmetic in the uncertain coordinates.

Let $x_i\in\mathbb{R}^d$, for $i=1,\ldots,n$, denote the $n$ frozen training
inputs.  Let $k(x,x_i)$ denote the supported stationary covariance kernel, and
define the kernel vector
\[
k_x = \bigl(k(x,x_1),\ldots,k(x,x_n)\bigr)^T.
\]
Let $\alpha\in\mathbb{R}^n$ denote the frozen posterior coefficient vector,
$L\in\mathbb{R}^{n\times n}$ the frozen nonsingular lower-triangular Cholesky
factor, and $m\in\mathbb{R}$ the frozen scalar mean.  The internal latent
posterior mean is
\[
\mu_{\mathrm{int}}(x)=m+k_x^T\alpha.
\]
Define $v(x)\in\mathbb{R}^n$ as the unique solution of
\[
L v(x)=k_x.
\]
The internal latent posterior variance is
\[
s^2_{\mathrm{int}}(x)=k(x,x)-v(x)^T v(x).
\]
All occurrences of the uncertain vector $x$ in these expressions are evaluated
with AutoIA validated values.  The frozen quantities $x_i$, $\alpha$, $L$, and
$m$ are treated as stored binary64 data.

If the fitted model uses an affine output standardization, let $a\in\mathbb{R}$
be its stored offset and let $b>0$ be its stored scale.  The user-coordinate
posterior is
\[
\mu(x)=a+b\,\mu_{\mathrm{int}}(x),\qquad
s^2(x)=b^2 s^2_{\mathrm{int}}(x).
\]

The Version 1.0.2 semantic recognizer supports scalar exact GPs with RBF or
Matern covariance ($\nu\in\{1/2,3/2,5/2\}$), an optional positive output-scale
kernel wrapper, zero or scalar constant mean, Gaussian observation noise, and
the documented affine input/output transforms.  Recognition is deliberately
fail-closed.  AutoIA does not claim a general validated Torch execution model.
The statistical fit, autograd, training convergence, and accelerator arithmetic
are not certified; the proof statement is conditional on the frozen stored
binary64 fitted state.  An independent 80-digit audit checks 504 posterior
mean/variance containment witnesses across the four supported kernel profiles.

\end{document}
